\documentclass[tikz, border=15pt]{standalone}
\usepackage[utf8]{inputenc}
\usepackage[T1]{fontenc}
\usepackage{roboto-mono}
\renewcommand*\familydefault{\ttdefault}

\usepackage{tikz}
\usetikzlibrary{arrows.meta, positioning, calc, decorations.pathmorphing}

\definecolor{stagecolor}{HTML}{1F6FEB}    % Accent Blue
\definecolor{hazardcolor}{HTML}{D29922}   % Amber / Orange
\definecolor{regcolor}{HTML}{7D8590}      % Muted Slate / Gray
\definecolor{boundarycolor}{HTML}{8C959F} % Dashed Line Gray
\definecolor{arrowgreen}{HTML}{2EA043}    % High-contrast green (GitHub green)
\definecolor{arrowpurple}{HTML}{8250DF}   % Vibrant Purple (GitHub Purple)

% Load diagram building utilities (e.g. \drawblock)
\IfFileExists{diagrams/util.tex}{\input{diagrams/util.tex}}{% TikZ Utilities for Hardware Architecture Diagrams
% Provides reusable styles and macros for drawing modules with stacked IO ports.

\usetikzlibrary{calc}

% Pin label style for ports
\tikzset{
    pinlabel/.style={font=\footnotesize\ttfamily\bfseries, text=stagecolor}
}

% Counters for port row stacking (declared once globally)
\newcount\nL
\newcount\nR
\newcount\nRows
\newcount\iL
\newcount\iR

% Macro: \drawblock
% Draws a rectangular hardware module with title and auto-stacked IO port columns.
% Input ports are stacked vertically on the left; output ports on the right.
% #1: Module identifier prefix (e.g. csr_decode)
% #2: Center X
% #3: Center Y
% #4: Width
% #5: Title (typeset with \texttt)
% #6: Left ports list  {port_name/color, ...}
% #7: Right ports list {port_name/color, ...}
% #8: Bottom ports list {port_name/x_offset/color, ...}
% #9: Custom height override (if empty, auto-computed from row count)
\newcommand{\drawblock}[9]{%
  % Count left and right ports to determine row count
  \nL=0
  \foreach \p in {#6} { \global\advance\nL by 1 }
  \nR=0
  \foreach \p in {#7} { \global\advance\nR by 1 }
  \ifnum\nL>\nR \nRows=\nL \else \nRows=\nR \fi
  \ifnum\nRows<1 \nRows=1 \fi
  % Geometry constants
  \pgfmathsetmacro{\rowH}{0.80}%
  \pgfmathsetmacro{\topH}{0.80}%
  \pgfmathsetmacro{\botH}{0.50}%
  \pgfmathsetmacro{\calcH}{\topH + \nRows * \rowH + \botH}%
  \def\overrideH{#9}%
  \ifx\overrideH\empty
    \pgfmathsetmacro{\H}{\calcH}%
  \else
    \pgfmathsetmacro{\H}{#9}%
  \fi
  \pgfmathsetmacro{\xL}{#2 - (#4)/2}%
  \pgfmathsetmacro{\xR}{#2 + (#4)/2}%
  \pgfmathsetmacro{\yT}{#3 + \H/2}%
  \pgfmathsetmacro{\yB}{#3 - \H/2}%
  % Module box
  \draw[draw=stagecolor, line width=1.5pt, fill=stagecolor, fill opacity=0.15]
    (\xL, \yT) rectangle (\xR, \yB);
  % Title
  \node[anchor=north, font=\bfseries\normalsize, text=stagecolor]
    at (#2, \yT - 0.15) {\texttt{#5}};
  % Left side ports (stacked column, top-to-bottom)
  \iL=0
  \foreach \pname/\pcol in {#6} {
    \global\advance\iL by 1
    \pgfmathsetmacro{\py}{\yT - \topH - (\the\iL - 0.5)*\rowH}%
    \coordinate (#1-\pname) at (\xL, \py);
    \node[pinlabel, anchor=west, text=\pcol] at (\xL + 0.12, \py) {\detokenize\expandafter{\pname}};
  }
  % Right side ports (stacked column, top-to-bottom)
  \iR=0
  \foreach \pname/\pcol in {#7} {
    \global\advance\iR by 1
    \pgfmathsetmacro{\py}{\yT - \topH - (\the\iR - 0.5)*\rowH}%
    \coordinate (#1-\pname) at (\xR, \py);
    \node[pinlabel, anchor=east, text=\pcol] at (\xR - 0.12, \py) {\detokenize\expandafter{\pname}};
  }
  % Bottom ports
  \foreach \pname/\pxoff/\pcol in {#8} {
    \pgfmathsetmacro{\px}{#2 + \pxoff}%
    \coordinate (#1-\pname) at (\px, \yB);
    \node[pinlabel, anchor=south, text=\pcol] at (\px, \yB + 0.10) {\detokenize\expandafter{\pname}};
  }
}
}

\begin{document}
\begin{tikzpicture}[
    >=Latex,
    font=\small,
    csrunitbox/.style={
        draw=stagecolor, line width=1.5pt,
        fill=stagecolor, fill opacity=0.15, text opacity=1.0,
        inner sep=0pt
    },
    stagetitle/.style={
        font=\bfseries\small, text=stagecolor, align=center
    },
    stageline/.style={dashed, draw=boundarycolor, line width=1.2pt},
    thickline/.style={draw=arrowgreen, line width=2.6pt, line cap=butt, line join=miter},
    hazardline/.style={draw=hazardcolor, line width=1.4pt, line cap=butt, line join=miter}
]

% ========================================================
% LAYOUT NOTES
%   y =  1.40 : csr_request bus to WB (over csr_data, jogs into csr_wb row 1)
%   y =  0.60 : csr_decode row 1 (csr_request out)
%   y = -0.20 : csr_data (csr_data row 1 -> csr_wb row 2)
%   y = -4.40 : write_intent bus feeding the hazard unit (pipelined)
%   y = -5.30 : csr_wb_request write-back bus (passes below the pipeline
%               registers since it is not pipelined; hops over write_intent)
% ========================================================

% ========================================================
% STAGE TITLES & BOUNDARY DASHED LINES
% ========================================================
\node[stagetitle] at (6.00,  3.9) {Instruction Decode (ID)};
\node[stagetitle] at (15.63, 3.9) {Execute (EX)};
\node[stagetitle] at (19.45, 3.9) {Memory (MEM)};
\node[stagetitle] at (25.70, 3.9) {Write-Back (WB)};

\draw[stageline] (13.75, 4.4) -- (13.75, 3.2);
\draw[stageline] (17.50, 4.4) -- (17.50, 3.2);
\draw[stageline] (21.40, 4.4) -- (21.40, 3.2);

% ========================================================
% STAGES & PIPELINE REGISTERS
% ========================================================

% Decode Stage outline (cut off at bottom)
\draw[draw=stagecolor, line width=1.8pt, fill=none]
  (-1.0, 3.2) -- (13.0, 3.2) -- (13.0, -3.2)
  decorate[decoration={zigzag, segment length=12pt, amplitude=3pt}] { -- (-1.0, -3.2) }
  -- cycle;
\node[anchor=north west, font=\bfseries\normalsize, text=stagecolor] at (-0.7, 3.05) {\texttt{decode\_stage}};

% Pipeline registers (ID/EX, EX/MEM, MEM/WB); extended down through the
% write_intent bus to show it is pipelined alongside the CSR signals
\foreach \xl in {13.40, 17.15, 21.05} {
  \pgfmathsetmacro{\xr}{\xl + 0.70}
  \path[fill=regcolor, fill opacity=0.15]
    (\xl, 3.2) -- (\xr, 3.2) -- (\xr, -4.8)
    decorate[decoration={zigzag, segment length=8pt, amplitude=2.5pt}] { -- (\xl, -4.8) }
    -- cycle;
  \draw[draw=regcolor, line width=1.3pt]
    (\xl, 3.2) -- (\xr, 3.2) -- (\xr, -4.8)
    decorate[decoration={zigzag, segment length=8pt, amplitude=2.5pt}] { -- (\xl, -4.8) }
    -- cycle;
}

% Write-Back Stage outline (cut off at bottom)
\draw[draw=stagecolor, line width=1.8pt, fill=none]
  (23.10, 3.2) -- (29.70, 3.2) -- (29.70, -3.2)
  decorate[decoration={zigzag, segment length=12pt, amplitude=3pt}] { -- (23.10, -3.2) }
  -- cycle;
\node[anchor=north west, font=\bfseries\normalsize, text=stagecolor] at (23.40, 3.05) {\texttt{write\_back\_stage}};

% CONTAINER LAYER: csr_unit box
% 0.4 padding: decode_stage -> csr_unit -> csr_decode / csr_data
\draw[csrunitbox] (0.40, 2.20) rectangle (12.60, -2.80);
\node[anchor=south east, font=\bfseries\normalsize, text=stagecolor] at (12.45, -2.70) {\texttt{csr\_unit}};

% ========================================================
% MODULE BLOCKS WITH STACKED IO PORTS (via util.tex)
% ========================================================

% csr_decode: x in [0.80, 5.20], y in [1.80, -2.40], row 1 at y = 0.60
\drawblock{csr_decode}{3.00}{-0.30}{4.4}{csr\_decode}%
  {instruction/arrowgreen}%
  {csr_request/arrowgreen}%
  {write_intent/0.0/hazardcolor}%
  {4.20}

% csr_data: x in [6.90, 12.20], top at y = 1.00 (below the WB csr_request bus),
% rows at y = -0.20 and y = -1.00
\drawblock{csr_data}{9.55}{-0.45}{5.3}{csr\_data}%
  {csr_request/arrowgreen, csr_wb_request/arrowgreen}%
  {csr_data/arrowgreen}%
  {}%
  {}

% csr_wb: x in [23.60, 28.80], top at y = 1.80, rows at y = 0.60 and y = -0.20
\drawblock{csr_wb}{26.20}{0.35}{5.2}{csr\_wb}%
  {csr_request/arrowgreen, csr_data/arrowgreen}%
  {csr_wb_request/arrowgreen}%
  {}%
  {}

% ========================================================
% INTERCONNECT WIRING
% ========================================================

% 1. "instruction" input to csr_decode
% (driven from decode_stage internals, so the wire starts inside the stage)
\draw[thickline] (-0.60, 0.60) -- (csr_decode-instruction);

% 2. "csr_request": out of csr_decode to a split point; one arm drops into
%    csr_data row 1, the other rises over csr_data and runs through the
%    pipeline, jogging down into csr_wb row 1
\draw[thickline] (csr_decode-csr_request) -- (5.70, 0.60) |- (csr_data-csr_request);
\draw[thickline] (5.70, 0.60) -- (5.70, 1.40) -- (22.45, 1.40) |- (csr_wb-csr_request);
\fill[arrowgreen] (5.70, 0.60) circle (3.2pt);

% 3. "csr_data": csr_data row 1 straight through the pipeline into csr_wb row 2
\draw[thickline] (csr_data-csr_data) -- (csr_wb-csr_data);

% 4. "write_intent" hazard bus: csr_decode bottom -> down -> right under the
%    pipeline, with per-stage drops (_e, _m, _w) into the hazard unit
\draw[hazardline] (csr_decode-write_intent) -- (3.00, -4.40) -- (23.35, -4.40) -- (23.35, -7.00);
\foreach \xd/\lbl in {15.55/e, 19.45/m} {
  \draw[hazardline] (\xd, -4.40) -- (\xd, -7.00);
  \fill[hazardcolor] (\xd, -4.40) circle (2.2pt);
  \node[anchor=west, font=\footnotesize\bfseries\ttfamily, text=hazardcolor]
    at (\xd + 0.12, -6.15) {write\_intent\_\lbl};
}
\node[anchor=west, font=\footnotesize\bfseries\ttfamily, text=hazardcolor]
  at (23.47, -6.15) {write\_intent\_w};

% 5. "csr_wb_request" write-back bus: csr_wb right -> down below the pipeline
%    registers -> left -> up into csr_data left row 2 (not pipelined).
%    Hops (radius \hop) over each write_intent wire it crosses.
\def\hop{0.22}
\draw[thickline]
  (csr_wb-csr_wb_request) -- (29.25, 0.60) -- (29.25, -5.30)
  \foreach \xc in {23.35, 19.45, 15.55} {
    -- (\xc + \hop, -5.30) arc[start angle=0, end angle=180, radius=\hop]
  }
  -- (6.30, -5.30)
  -- (6.30, -4.40 - \hop) arc[start angle=-90, end angle=90, radius=\hop]
  |- (csr_data-csr_wb_request);

% ========================================================
% HAZARD UNIT (cut off on the bottom with sawtooth)
% ========================================================
\path[fill=hazardcolor, fill opacity=0.15]
  (-1.0, -7.00) -- (29.70, -7.00) -- (29.70, -8.20)
  decorate[decoration={zigzag, segment length=12pt, amplitude=3pt}] { -- (-1.0, -8.20) }
  -- cycle;
\draw[draw=hazardcolor, line width=1.5pt]
  (-1.0, -7.00) -- (29.70, -7.00) -- (29.70, -8.20)
  decorate[decoration={zigzag, segment length=12pt, amplitude=3pt}] { -- (-1.0, -8.20) }
  -- cycle;
\node[anchor=center, font=\bfseries\normalsize, text=hazardcolor] at (14.35, -7.60) {\texttt{hazard\_unit}};

\end{tikzpicture}
\end{document}
